\par\Needspace{6\baselineskip}\medskip
\hypertarget{native-125-article092}{}\pdfbookmark[2]{92.}{article-092}
\begin{center}92.\end{center}
Quamquam aequationes 15, $15^{\ast}$ in quibusdam casibus tres radices reales habere possint, tamen ambiguum nunquam erit, quamnam in problemate nostro adoptare oporteat. Quum enim $h$ manifesto sit quantitas positiva, ex aequationum theoria facile concluditur, aequationem 15 habere radicem unicam positivam vel cum duabus imaginariis vel cum duabus negativis: iam quum $y=\dfrac{m}{\BareRoot{(l+x)}}$ necessario esse debeat quantitas positiva, nullam hic incertitudinem \SourceLeaf{126}remanere patet. Quod vero attinet ad aequationem $15^{\ast}$, primo observamus, $L$ necessario esse maiorem quam 1: quod facile probatur, si aequatio in art. 89 tradita sub formam $L=1+\dfrac{\cos\tfrac12 f^2}{-\cos f}+\dfrac{\operatorname{tang}2\omega^2}{-\cos f}$ ponitur. Porro substituendo in aequatione $12^{\ast}$ pro $M$, $Y\BareRoot{(L-x)}$, prodit $Y+1=(L-x)X$, adeoque
\[
\begin{aligned}
Y+1>(1-x)X>{}&\tfrac43+\dfrac{4}{3\mathbin{.}5}x+\dfrac{4\mathbin{.}6}{3\mathbin{.}5\mathbin{.}7}xx
\\[4pt]
&+\dfrac{4\mathbin{.}6\mathbin{.}8}{3\mathbin{.}5\mathbin{.}7\mathbin{.}9}x^3+\text{etc.}>\tfrac43,
\end{aligned}
\]
et proin $Y>\tfrac13$. Statuendo itaque $Y=\tfrac13+Y'$, necessario $Y'$ erit quantitas positiva, aequatio $15^{\ast}$ autem hinc transit in hanc
\[
Y'^3+2Y'Y'+(1-H)Y'+\tfrac{4}{27}-\tfrac29H=0,
\]
quam plures radices positivas habere non posse ex aequationum theoria facile probatur. Hinc colligitur, aequationem $15^{\ast}$ unicam radicem habituram esse maiorem quam $\tfrac13$\footnote{Siquidem problema revera solubile esse supponimus.}, quam neglectis reliquis in problemate nostro adoptare oportebit\HistoricalNoteDouble{Handschriftliche Bemerkung:] Die Gleichung $15^{\ast}$ hat
\par 1) Eine reelle negative Wurzel, wenn $H$ zwischen den Grenzen 0 und $\dfrac{+1+\BareRoot{5}}{6}$ liegt, nebst zwei imaginären.
\par 2) Drei reelle Wurzeln, worunter Eine positive, wenn $H$ zwischen $\dfrac{+1-\BareRoot{5}}{6}$ und 0.
\par 3) Eine reelle positive Wurzel und zwei imaginäre, wenn $H$ zwischen $-\infty$ und $\dfrac{+1-\BareRoot{5}}{6}\mathpunct{\cdot}$
\par 4) Drei reelle Wurzeln, unter denen Eine negativ, wenn $H$ zwischen $\dfrac{+1+\BareRoot{5}}{6}$ und $+\infty$.
\par Offenbar kann also nur von Fall 4 hier die Rede sein, wo sich zwei positive Wurzeln finden. Allein die Eine derselben ist hier immer kleiner als $\dfrac{-1+\BareRoot{5}}{6}\mathpunct{\raisebox{.45ex}{,}}$ die andere grösser; letztere kann also allein gültig sein.
\par Bei dem im Text angezeichneten Verfahren ist klar, dass die letzte Gleichung nur dann zwei positive Wurzeln haben könnte, wenn zugleich $1-H$ negativ und $\tfrac{4}{27}-\tfrac29H$ positiv wäre, welches offenbar unmöglich ist, da $\tfrac{4}{27}-\tfrac29H=\tfrac29(1-H)-\tfrac{2}{27}$.
}.
